\documentclass[12pt, a4paper]{article}

% ── Fonts & Encoding ──────────────────────────────────────────────────────────
\usepackage{fontspec}
\usepackage{xeCJK}
\setCJKmainfont{DFKai-SB}
\setmainfont{Times New Roman}

% ── Layout ────────────────────────────────────────────────────────────────────
\usepackage[top=2.5cm, bottom=2.5cm, left=3cm, right=3cm]{geometry}
\usepackage{array, booktabs, tabularx, multirow, makecell}
\usepackage{enumitem}
\usepackage{xcolor}
\usepackage{mdframed}
\usepackage{titlesec}
\usepackage{parskip}
\usepackage{microtype}
\usepackage{amssymb}
\usepackage{hyperref}

% ── Colours ───────────────────────────────────────────────────────────────────
\definecolor{prismblue}{RGB}{30, 90, 160}
\definecolor{defbox}{RGB}{240, 246, 255}
\definecolor{defborder}{RGB}{100, 150, 220}
\definecolor{primebox}{RGB}{255, 248, 230}
\definecolor{primeborder}{RGB}{200, 160, 60}
\definecolor{lightgray}{RGB}{245,245,245}

% ── Section style ─────────────────────────────────────────────────────────────
\titleformat{\section}{\large\bfseries\color{prismblue}}{}{0em}{}[\titlerule]
\titleformat{\subsection}{\normalsize\bfseries}{}{0em}{}

% ── Likert helper ─────────────────────────────────────────────────────────────
\newcommand{\likertfive}[2]{%
  \begin{tabular*}{\linewidth}{@{\extracolsep{\fill}}ccccc@{}}
    \footnotesize 1 & \footnotesize 2 & \footnotesize 3 & \footnotesize 4 & \footnotesize 5 \\[2pt]
    $\bigcirc$ & $\bigcirc$ & $\bigcirc$ & $\bigcirc$ & $\bigcirc$
  \end{tabular*}\par\vspace{3pt}%
  \begin{tabular*}{\linewidth}{@{}l@{\extracolsep{\fill}}r@{}}
    \footnotesize #1 & \footnotesize #2
  \end{tabular*}%
}

% ── Choice helper ─────────────────────────────────────────────────────────────
\newcommand{\choiceitem}[1]{\item[$\square$] #1}
\newcommand{\radioitem}[1]{\item[$\bigcirc$] #1}

% ── Question counter ──────────────────────────────────────────────────────────
\newcounter{qnum}
\newcommand{\q}[1]{\stepcounter{qnum}\noindent\textbf{Q\arabic{qnum}.}\enspace #1\par\smallskip}

% ─────────────────────────────────────────────────────────────────────────────
\begin{document}

% ── Cover ─────────────────────────────────────────────────────────────────────
\begin{center}
  {\LARGE\bfseries\color{prismblue} 漫畫消費與盜版認知調查}\\[0.4em]
  {\large Prism 政策感知研究計畫}\\[0.3em]
  {\small 本問卷針對台灣大學生，預計填寫時間約 8--10 分鐘。\\
   所有回答均用於學術研究，匿名處理，不作商業用途。}
\end{center}

\bigskip

% ── Moral Priming Block ───────────────────────────────────────────────────────
\begin{mdframed}[backgroundcolor=primebox, linecolor=primeborder,
                 linewidth=1.2pt, innerleftmargin=12pt, innerrightmargin=12pt,
                 innertopmargin=10pt, innerbottommargin=10pt]
\begin{center}
  \textbf{$\bigstar$ 回顧這一年的漫畫時光 $\bigstar$}
\end{center}
\smallskip
這一年，你讀過幾部漫畫？在每個深夜追更新、每次被劇情震撼的瞬間，背後都有作者、編輯和整個製作團隊的心血在支撐著這個故事。

每一個閱讀選擇，都在創作生態中留下了不同的痕跡。

\medskip
\textit{（以下問題請根據你的真實經驗作答，沒有對錯之分）}
\end{mdframed}

\bigskip

% ═══════════════════════════════════════════════════════════════
\section*{A 部分｜閱讀習慣}
% ═══════════════════════════════════════════════════════════════

\q{你平均多久閱讀一次漫畫（含紙本、電子、任何形式）？}
\begin{itemize}[leftmargin=2em, itemsep=2pt]
  \radioitem{幾乎每天}
  \radioitem{每週數次}
  \radioitem{每週約一次}
  \radioitem{每月數次}
  \radioitem{很少或幾乎不讀}
\end{itemize}

\bigskip

\q{你主要透過哪些管道閱讀漫畫？（可複選）}
\begin{itemize}[leftmargin=2em, itemsep=2pt]
  \choiceitem{官方付費／訂閱平台（如 LINE Manga、Comico、ComiXology、BookWalker 等）}
  \choiceitem{出版社官方網站或 APP}
  \choiceitem{圖書館借閱或購買實體書}
  \choiceitem{非官方免費閱讀網站（如各類漫畫閱讀網）}
  \choiceitem{社群媒體（Instagram、Twitter/X、Threads、Facebook 等）}
  \choiceitem{通訊軟體（LINE 群組、Discord 等）}
  \choiceitem{朋友直接傳送給我}
  \choiceitem{其他：\underline{\hspace{6cm}}}
\end{itemize}

\bigskip

\q{過去一年，你是否曾透過「非官方」管道閱讀漫畫？}
\begin{itemize}[leftmargin=2em, itemsep=2pt]
  \radioitem{是}
  \radioitem{否}
  \radioitem{不確定（不清楚所使用的管道是否為官方授權）}
\end{itemize}

\bigskip

\q{你曾經做過以下哪些行為？（可複選）}
\begin{itemize}[leftmargin=2em, itemsep=2pt]
  \choiceitem{在非官方網站上瀏覽漫畫}
  \choiceitem{將漫畫頁面下載或儲存至個人裝置}
  \choiceitem{在社群媒體上分享漫畫截圖（1至數格）}
  \choiceitem{將多頁或整章漫畫傳送給朋友}
  \choiceitem{在論壇或群組中上傳漫畫頁面}
  \choiceitem{以上皆無}
\end{itemize}

% ═══════════════════════════════════════════════════════════════
\section*{B 部分｜認知與意識}
% ═══════════════════════════════════════════════════════════════

\q{當你在線上閱讀漫畫時，你多常留意該來源是否為官方授權？}

\medskip
\likertfive{從不留意}{總是留意}

\bigskip

\q{你是否能分辨官方平台與非官方來源？}
\begin{itemize}[leftmargin=2em, itemsep=2pt]
  \radioitem{能清楚分辨}
  \radioitem{大多能分辨}
  \radioitem{有時難以分辨}
  \radioitem{通常無法分辨}
\end{itemize}

\bigskip

\q{若你曾使用非官方來源閱讀漫畫，主要原因為何？（可複選）}
\begin{itemize}[leftmargin=2em, itemsep=2pt]
  \choiceitem{免費，不需付費}
  \choiceitem{更新速度比官方更快}
  \choiceitem{官方平台沒有我想看的作品（如未代理、已下架）}
  \choiceitem{習慣，沒有特別思考來源}
  \choiceitem{朋友推薦或傳送給我}
  \choiceitem{不知道有對應的官方管道}
  \choiceitem{我未曾使用非官方來源}
\end{itemize}

% ═══════════════════════════════════════════════════════════════
\section*{C 部分｜定義確認}
% ═══════════════════════════════════════════════════════════════

\begin{mdframed}[backgroundcolor=defbox, linecolor=defborder,
                 linewidth=1.2pt, innerleftmargin=14pt, innerrightmargin=14pt,
                 innertopmargin=10pt, innerbottommargin=10pt]
\textbf{本問卷對「盜版漫畫」的定義如下：}

\medskip
\textit{任何未經作者或出版社授權而上傳、流通的漫畫版本，無論其形式為網站、論壇、社群媒體或通訊軟體，均屬盜版。}

\medskip
\footnotesize 補充說明：根據現行著作權法，「散布」及「重製」未授權內容屬於侵權行為。單純「瀏覽」在法律上的定性較為複雜，各方見解不一。
\end{mdframed}

\bigskip

\q{上述定義是否與你原本對「盜版」的理解一致？}
\begin{itemize}[leftmargin=2em, itemsep=2pt]
  \radioitem{完全一致}
  \radioitem{大致一致，但我的定義較寬鬆}
  \radioitem{大致一致，但我的定義較嚴格}
  \radioitem{差異很大}
\end{itemize}

\bigskip

\q{根據上述定義，你認為你過去的閱讀行為涉及盜版的程度為何？}

\medskip
\likertfive{完全沒有涉及}{幾乎都涉及}

% ═══════════════════════════════════════════════════════════════
\section*{D 部分｜態度與感受}
% ═══════════════════════════════════════════════════════════════

\textit{以下兩題請根據你讀完上述定義後的真實感受作答。}

\bigskip

\q{【個人感受】了解上述定義後，你對自己過去的閱讀習慣有什麼想法或感受？}

\medskip
\likertfive{完全沒有特別感受}{非常在意，想重新審視習慣}

\bigskip

\q{【社交感受】如果你的朋友或家人能看到你過去所有的漫畫閱讀紀錄，你會有什麼感受？}

\medskip
\likertfive{完全不在意}{非常在意或感到不自在}

\bigskip

\q{讀完上述定義後，你對「盜版」這件事的看法是否有所改變？}
\begin{itemize}[leftmargin=2em, itemsep=2pt]
  \radioitem{是，我認為盜版比我原本想的更嚴重}
  \radioitem{是，我認為盜版沒有我原本想的那麼嚴重}
  \radioitem{看法沒有改變}
  \radioitem{我本來就沒有使用盜版，此問不適用}
\end{itemize}

% ═══════════════════════════════════════════════════════════════
\section*{E 部分｜合理使用範圍}
% ═══════════════════════════════════════════════════════════════

\textit{請針對以下各情境，分別回答兩個問題：}
\begin{itemize}[noitemsep]
  \item \textbf{甲} — 你認為這個行為屬於盜版嗎？（1 = 完全不是，5 = 絕對是）
  \item \textbf{乙} — 你認為這個行為應該被允許嗎？（○ 應該 ／ ○ 不應該 ／ ○ 無所謂）
\end{itemize}

\bigskip

\setlength{\tabcolsep}{6pt}
\renewcommand{\arraystretch}{2.2}

\begin{tabular}{@{} p{5.0cm} p{5.5cm} p{3.5cm} @{}}
\toprule
\textbf{情境} & \textbf{甲\enspace 是否屬於盜版？} & \textbf{乙\enspace 應否被允許？} \\
\midrule

\makecell[l]{E1. 在社群媒體分享\\1–2 格漫畫截圖作為感想} &
\likertfive{不是}{是} &
$\bigcirc$ 應該 \newline $\bigcirc$ 不應該 \newline $\bigcirc$ 無所謂 \\

\makecell[l]{E2. 在論壇發文附上\\單集不超過 1/3 的頁面\\（如 PTT C\_Chat 慣例）} &
\likertfive{不是}{是} &
$\bigcirc$ 應該 \newline $\bigcirc$ 不應該 \newline $\bigcirc$ 無所謂 \\

\makecell[l]{E3. 將漫畫下載至個人裝置\\供離線閱讀，不對外分享} &
\likertfive{不是}{是} &
$\bigcirc$ 應該 \newline $\bigcirc$ 不應該 \newline $\bigcirc$ 無所謂 \\

\makecell[l]{E4. 在論壇或群組\\上傳完整單集漫畫} &
\likertfive{不是}{是} &
$\bigcirc$ 應該 \newline $\bigcirc$ 不應該 \newline $\bigcirc$ 無所謂 \\

\makecell[l]{E5. 使用或架設提供\\完整系列的免費閱讀網站} &
\likertfive{不是}{是} &
$\bigcirc$ 應該 \newline $\bigcirc$ 不應該 \newline $\bigcirc$ 無所謂 \\

\bottomrule
\end{tabular}

\bigskip

\q{以你個人的標準，你認為「合理使用」的界線在哪裡？有什麼原則讓某些行為在你看來是可以接受的？}

\vspace{3cm}
\hrule
\vspace{1cm}
\hrule

\bigskip

% ═══════════════════════════════════════════════════════════════
\section*{F 部分｜推廣策略感知}
% ═══════════════════════════════════════════════════════════════

\textit{本部分與版權定義無關，旨在了解你對不同推廣訊息的感受與反應。}

\bigskip

\q{你是否曾在任何平台（音樂、影音、漫畫等）看過「年度回顧」功能，例如 Spotify Wrapped？}
\begin{itemize}[leftmargin=2em, itemsep=2pt]
  \radioitem{有，而且很喜歡看}
  \radioitem{有，但沒什麼特別感覺}
  \radioitem{有看過，但不太在意}
  \radioitem{沒有印象}
\end{itemize}

\bigskip

\q{如果漫畫官方平台在年底推出功能，顯示：「今年你支持了 12 位創作者，為創作生態注入了活力」，你的感受為何？}

\medskip
\likertfive{完全沒有共鳴}{非常有共鳴，覺得很有意義}

\bigskip

\q{上述這類訊息是否會影響你選擇閱讀管道的決定？}
\begin{itemize}[leftmargin=2em, itemsep=2pt]
  \radioitem{會，我會更傾向使用官方平台}
  \radioitem{可能有一點影響}
  \radioitem{不會，不影響我的選擇}
  \radioitem{反而覺得這種訊息有點刻意或說教}
\end{itemize}

\bigskip

\q{以下四種訊息風格，哪一種你覺得最有共鳴？（單選）}
\begin{itemize}[leftmargin=2em, itemsep=2pt]
  \radioitem{「今年你支持了 X 位創作者，讓他們的故事得以繼續。」（情感連結）}
  \radioitem{「今年你透過官方管道閱讀了 X 部作品。」（中性事實）}
  \radioitem{「如果你今年的閱讀都透過官方管道，你最愛的作者將多獲得 NT\$X 的收入。」（反事實影響）}
  \radioitem{「你和 X\% 的讀者一樣，選擇用支持的方式享受漫畫。」（社會規範）}
\end{itemize}

\bigskip

\q{哪一種訊息風格你覺得最不自然或最像在說教？（單選）}
\begin{itemize}[leftmargin=2em, itemsep=2pt]
  \radioitem{情感連結型（「讓他們的故事得以繼續」）}
  \radioitem{中性事實型（純數字統計）}
  \radioitem{反事實影響型（「你的選擇讓作者多賺了 NT\$X」）}
  \radioitem{社會規範型（「X\% 的讀者選擇支持」）}
\end{itemize}

\bigskip

\q{你認為什麼樣的推廣方式，最可能讓你考慮改用或繼續使用官方閱讀管道？（開放題）}

\vspace{3cm}
\hrule
\vspace{1cm}
\hrule

\bigskip

\begin{center}
  \textit{感謝你的參與。你的回答將協助我們了解消費者視角，\\
  並為漫畫產業的版權政策提供更貼近現實的參考依據。}
\end{center}

% ═══════════════════════════════════════════════════════════════
% APPENDIX — AI Simulation Persona Prompt
% ═══════════════════════════════════════════════════════════════
\newpage
\section*{附錄：AI 模擬人設提示詞（Persona Prompt）}
\textit{供 Prism 引擎模擬用，不隨問卷印出。}

\bigskip

\begin{mdframed}[backgroundcolor=lightgray, linewidth=0.5pt,
                 innerleftmargin=14pt, innerrightmargin=14pt]
\textbf{〔人設底層提示詞 — 繁體中文版〕}

\medskip
你是一位台灣的大學生，年齡介於 18 至 25 歲之間。\texttt{[segment\_description]}

你平時有閱讀漫畫的習慣，對於不同的閱讀平台有一定的使用經驗。
請以第一人稱誠實回答以下問卷，根據你自身的習慣、想法與感受作答。
回答時不需要顯得「政治正確」或「道德高尚」，請反映你真實的行為與態度。

\medskip
\textbf{你的基本背景設定：}
\begin{itemize}[noitemsep]
  \item 閱讀頻率：\texttt{[frequency]}
  \item 慣用平台類型：\texttt{[platform\_types]}
  \item 對版權議題的關注程度：\texttt{[copyright\_awareness\_level]}（低 / 中 / 高）
  \item 是否曾使用非官方來源：\texttt{[has\_used\_piracy]}（是 / 否 / 不確定）
  \item 社交圈對盜版的態度：\texttt{[peer\_norm]}（普遍接受 / 中立 / 普遍反對）
\end{itemize}

\medskip
\textbf{回答格式說明：}
\begin{itemize}[noitemsep]
  \item 單選題：回答選項編號（如：\texttt{B}）
  \item 複選題：回答所有選擇的編號（如：\texttt{A, C, E}）
  \item Likert 題：回答 1 至 5 的整數
  \item 開放題：以 1 至 3 句話作答，使用自然、口語的繁體中文
\end{itemize}
\end{mdframed}

\bigskip

\textbf{段別（Segment）範例設定：}

\medskip
\begin{tabularx}{\textwidth}{@{} l X X X X @{}}
\toprule
\textbf{段別} & \textbf{描述} & \textbf{閱讀頻率} & \textbf{慣用平台} & \textbf{盜版使用} \\
\midrule
重度盜版用戶 & 主要依賴非官方來源，認為盜版無害 & 每天 & 非官方網站 & 是 \\
混合型用戶 & 官方非官方皆用，依作品決定 & 每週數次 & 混合 & 是 \\
官方付費用戶 & 習慣付費，偶爾接受朋友分享 & 每週數次 & 官方平台 & 不確定 \\
輕度閱讀者 & 偶爾看，不特別在意來源 & 每月數次 & 社群媒體 & 不確定 \\
不閱讀者 & 幾乎不讀漫畫，無相關習慣 & 幾乎不讀 & 無 & 否 \\
\bottomrule
\end{tabularx}

\end{document}
